% ====================================================================
% 第7章：用户痛点（LaTeX版本，含TikZ图）
% ====================================================================

\section{用户痛点}\label{ch:painpoints}

\subsection{研究方法}

本章基于\textbf{The Mom Test}方法论，从Reddit (62帖)、Hacker News (46帖)和X/Twitter (23帖)的真实社区讨论中提取用户痛点。遵循以下原则：(1) 从实际用户体验提取，非理论观点；(2) 关注用户做了什么，而非声称想要什么；(3) 识别变通方案作为未满足需求的证据。

\subsection{痛点分布}

共提取\textbf{97个独立痛点}，分为15个类别。

\begin{figure}[htbp]
\centering
\begin{tikzpicture}[scale=0.85]
\begin{axis}[
  xbar,
  xmin=0, xmax=15,
  y=7mm,
  bar width=5mm,
  xlabel={痛点数量},
  xlabel style={font=\small},
  ytick=data,
  yticklabels={
    {生产环境部署鸿沟},
    {Agent可靠性天花板},
    {自改进可行性存疑},
    {框架/工具缺陷},
    {知识与记忆持久化},
    {评估与基准失真},
    {安全/成本/治理},
    {人-Agent协作},
    {定义与标准缺失},
    {恶意进化风险},
    {交互模式错配},
    {从业者可及性},
    {绝对增益有限},
    {企业治理},
    {生产证明缺失}
  },
  yticklabel style={font=\scriptsize, anchor=east},
  nodes near coords,
  nodes near coords style={font=\tiny},
  every axis plot/.append style={fill=cat1!60, draw=cat1},
  grid=major,
  grid style={gray!20},
  width=12cm,
  height=14cm,
  title={97个用户痛点分类分布},
  title style={font=\small\bfseries},
]
\addplot coordinates {
  (12,14) (12,13) (11,12) (10,11) (9,10) (7,9) (5,8) (5,7) (3,6) (3,5)
  (3,4) (2,3) (2,2) (1,1) (1,0)
};
\end{axis}
\end{tikzpicture}
\caption{Mom Test用户痛点分布（N=97）。数据来源于Reddit、HN和X/Twitter的真实社区讨论。}
\label{fig:ch7-painpoints}
\end{figure}

\subsection{跨平台主题}

\subsubsection{Theme 1: ``Demo有效，生产失败''}

\begin{itemize}[nosep]
\item Reddit: ``80\% reliability ceiling'' — agents通过benchmark但在真实workflow中失败
\item HN: ``Nobody can demonstrate real production success with self-improving agents''
\item X: 改进百分比有限（DGM: 20\%$\to$50\% on SWE-bench）
\end{itemize}

\subsubsection{Theme 2: ``自改进是神话（目前）''}

\begin{itemize}[nosep]
\item Reddit: ``Feedback loops require human labor, drift, plateau without verifiers''
\item HN: ``LLMs are bad at prompting other LLMs, the prompt search space is too large''
\item GitHub Issue: ``Self-improving AI agent is a myth''
\end{itemize}

\subsubsection{Theme 3: ``框架被抛弃''}

\begin{itemize}[nosep]
\item Reddit: ``Zero prompt visibility, 1GB dependency bloat, state management nightmares''
\item HN: ``LangChain abstractions hide debugging, frameworks fail in production''
\item 开发者: ``AutoGen, LangGraph, CrewAI, PydanticAI, Swarm, Agno — zero survive months of real use''
\end{itemize}

\subsubsection{Theme 4: ``Benchmark被操控''}

\begin{itemize}[nosep]
\item Reddit: ``No session-level evals, benchmarks ignore 92\% of labor''
\item HN: ``SWE-Bench is contaminated/saturated/gameable, Goodhart's Law applies to every metric''
\item X: ``Research optimizes for benchmarks, not real-world value''
\end{itemize}

\subsubsection{Theme 5: ``自修改是危险的''}

\begin{itemize}[nosep]
\item Reddit: ``RCE vulnerabilities in open platforms, agents burning \$2K overnight''
\item HN: ``Self-modification is a prompt injection attack vector. The skills hub is poisoned.''
\item X: ``Agents evolve in harmful, unintended directions''
\end{itemize}

\subsection{痛点优先级矩阵}

\begin{figure}[htbp]
\centering
\resizebox{\textwidth}{!}{%
\begin{tabular}{cl
  >{\centering\arraybackslash}p{14mm}
  >{\centering\arraybackslash}p{14mm}
  >{\centering\arraybackslash}p{40mm}}
\toprule
\textbf{优先级} & \textbf{痛点} & \textbf{频率} & \textbf{严重度} & \textbf{当前解决方案} \\
\midrule
1 & 无持久记忆/跨会话学习 & 极高 & 致命 & ReasoningBank, SEDM（研究阶段） \\
2 & 自进化评估成本过高 & 高 & 高 & 无实用方案 \\
3 & 自主agent成本透明度 & 高 & 致命 & 仅手动监控 \\
4 & Agent自修改无回滚/版本控制 & 中 & 高 & 无 \\
5 & 框架生产环境不稳定 & 高 & 中 & 多个竞争选项 \\
6 & 安全/沙箱逃逸担忧 & 中 & 高 & Docker/VM（不足） \\
7 & 静态智能—无实时学习 & 高 & 根本性 & 微调（仅离线） \\
8 & 进化方法"扔意大利面"观感 & 中 & 中 & AlphaEvolve（未公开） \\
\bottomrule
\end{tabular}%
}
\caption{痛点优先级矩阵。按频率和严重度排序，列出当前最佳解决方案状态。}
\label{fig:ch7-priority}
\end{figure}



\begin{table}[htbp]
\centering
\small
\begin{tabular}{p{30mm}p{34mm}p{38mm}p{38mm}}
\toprule
\textbf{平台} & \textbf{主导痛点} & \textbf{典型证据形态} & \textbf{对评估/产品的含义} \\
\midrule
Reddit & 生产可靠性、成本失控、框架调试困难 & 长线程复盘、失败案例、替代工具比较 & 需要会话级回归测试、预算上限和可观测性 \\
Hacker News & 自改进可证伪性、benchmark污染、工程边界 & 反例论证、系统设计讨论、历史类比 & 需要独立验证集、负结果记录和安全门控 \\
X/Twitter & 短期增益、demo传播、热点系统比较 & 数字化增益摘录、论文/产品链接、观点传播 & 需要把传播指标与可复现实验分离 \\
GitHub Issues & 安装、依赖、权限、长期维护 & bug报告、feature request、维护者回复 & 需要最小可运行样例、版本锁定和回滚路径 \\
\bottomrule
\end{tabular}
\caption{跨平台痛点证据类型对照。不同平台提供互补信号：Reddit/HN更适合发现生产与方法论风险，X/Twitter更适合捕捉传播叙事，GitHub Issues更适合验证工程落地摩擦。}
\label{tab:ch7-platform-evidence}
\end{table}

\subsection{对awesome-evolution项目的启示}

基于这些发现，最有价值的资源方向是：

\begin{enumerate}[nosep]
\item \textbf{生产就绪框架}：具有已验证部署记录的系统
\item \textbf{评估工具包}：超越benchmark的session级评估
\item \textbf{安全/治理模式}：针对自修改agent的安全沙箱和回滚机制
\item \textbf{记忆架构}：不膨胀上下文的持久记忆方案
\item \textbf{实用自改进模式}：无需人工验证循环的可行改进路径
\end{enumerate}

% Expanded-file safety copy. Master instructed all Ch7 work to move from
% ch7-painpoints.tex to ch7-painpoints-expanded.tex after repeated overwrite
% conflicts. This file preserves the currently available Ch7 content and is
% the only Ch7 target this session will touch going forward.
